\section{Averaging matrix lower bounds}
\label{sec:matrix_floor_pin}

\begin{theorem}[Averaging two lower bounds]
\label{thm:matrix_floor_pin}
\lean{Matrix.floor_pin_le}
\leanok
Let $A,P$ be complex matrices on a finite-dimensional space and let
$b,c\in\mathbb R$. If $2bI\le A$ and $cP\le A$, then
\begin{align}
 b(I-P)+\left(b+\frac c2\right)P\le A.
\end{align}
No projection or commutation assumption on $P$ is needed. This is the
averaging step preceding the logarithmic comparison in
\cite[Section~7, lines~603--615]{OpenAI2026AreaLaw}.
\end{theorem}
\begin{proof}
\leanok
Add the two assumed inequalities and divide by two. The resulting lower
matrix $bI+(c/2)P$ equals the displayed expression.
\end{proof}

\begin{theorem}[Adding the complementary term]
\label{thm:matrix_identity_extension_floor}
\lean{Matrix.identity_extension_floor}
\leanok
Let $T,S$ be complex matrices on a finite-dimensional space and let
$c\in\mathbb R$. If $S\le I$, $c\le1$ and $cS\le T$, then
\begin{align}
 cI\le T+(I-S).
\end{align}
In particular, this applies to extension by the identity off the range
of an orthogonal projection. The inequality itself does not require
$S$ to be a projection. This auxiliary order inequality extends a lower
bound to the whole space.
\end{theorem}
\begin{proof}
\leanok
Since $I-S\ge0$ and $c\le1$, one has $c(I-S)\le I-S$.
Add this inequality to $cS\le T$.
\end{proof}
